\documentclass[tikz]{standalone}
\usepackage{../sty/phylogenetic-diagrams} % Include the style file containing the code defining the structure and style of the document

\begin{document}
\begin{tikzpicture}[node distance=2.5cm]
	\node (time_archetype) at (-3,0) {$t_{\alpha}$};
	\node (time_d1) at (-3,-2) {$t_{\delta_1}$};
	\node (time_d2) at (-3,-3) {$t_{\delta_2}$};
	\node (time_0) at (-3,-4) {};
	\node[hypotheticalChild] (archetype) at (-2,0) {};
	\node[hypotheticalChild] (birth_d1) at (-2,-2) {};
	\node[hypotheticalChild] (birth_d2) at (-2,-3) {};
	\node (archetype_lineage) at (-2,-4) {};
	\node (d1_lineage) at (0,-4) {};
	\node (d2_lineage) at (-1,-4) {};
	\draw[undirected] (archetype) -- (birth_d1) -- (birth_d2);
	\draw[undirected, dashed] (birth_d2) -- (archetype_lineage);
	\draw[undirected] (birth_d1) -| (d1_lineage);
	\draw[undirected] (birth_d2) -| (d2_lineage);
	\draw[undirected, dashed] (time_archetype) -- (time_d1) -- (time_d2) -- (time_0);
\end{tikzpicture}
\end{document}